% ============================================================
% CAPITOLO 9 — Progetto 6: Frontend Web con React e Autenticazione
% ============================================================
\chapter{Progetto 6: Frontend Web con React e Autenticazione}\index{React}\index{Tailwind CSS}\index{Vite}\index{SKILL.md}

\section*{Cosa Costruirai}
Un'applicazione frontend React completa che:
\begin{itemize}
  \item Si collega al backend Notes API
  \item Permette login con Google e GitHub
  \item Mostra una dashboard con le note dell'utente
  \item Supporta operazioni CRUD complete
  \item Ha rotte protette
  \item Usa un design system coerente con Tailwind~CSS
\end{itemize}

\textbf{Tempo stimato}: 60--90 minuti\\
\textbf{Prerequisito}: Backend Notes API con autenticazione (Capitoli~6--8)

\begin{notebox}
\textbf{Box Tooling --- Stack scelto per questo esempio.}
\begin{itemize}
  \item \textbf{Framework:} React~19 con Vite
  \item \textbf{Styling:} Tailwind~CSS~4
  \item \textbf{Routing:} React Router~7
  \item \textbf{HTTP Client:} Axios
\end{itemize}
\textbf{Alternative equivalenti:} Vue.js/Nuxt, Svelte/SvelteKit, Angular, Next.js, Astro. Il \textbf{pattern} (componenti, routing protetto, gestione stato, chiamate API) \`e lo stesso in qualunque framework frontend moderno. Il file \texttt{SKILL.md} che creerai in questo capitolo funziona con qualsiasi stack --- basta adattare le convenzioni.
\end{notebox}

% ---------------------------------------------------------
\section{Setup del Progetto Frontend}

\begin{praticabox}[title={Pratica --- Creare il progetto React}]
Apri un terminale nella directory padre del backend:
\begin{lstlisting}[language=bash]
cd ..
npm create vite@latest notes-frontend -- --template react
cd notes-frontend
npm install
npm install react-router-dom axios
npm install -D tailwindcss @tailwindcss/vite
\end{lstlisting}
Ora hai due progetti affiancati:
\begin{lstlisting}[language={}]
workspace/
+-- notes-api/          <- Backend (Capitoli 6-8)
+-- notes-frontend/     <- Frontend (questo capitolo)
\end{lstlisting}
\end{praticabox}

% ---------------------------------------------------------
\section{Agent Skill per il Frontend}

Questo è il momento di introdurre il tuo primo \textbf{SKILL.md} --- un documento che conferisce all'IA competenze specializzate.

\begin{praticabox}[title={Pratica --- Crea il file \texttt{SKILL.md}}]
\end{praticabox}

\begin{lstlisting}[language={},title={\texttt{SKILL.md} --- React Frontend Developer}]
# React Frontend Developer Skill

## Principi di Design
- Mobile-first
- Componenti piccoli: ogni componente fa UNA cosa
- Nessuna logica di business nei componenti UI
- Accessibilita': elementi interattivi con supporto tastiera

## Struttura Componenti
src/
  components/
    ui/         <- Button, Input, Card, Modal
    layout/     <- Header, Footer, PageContainer
    notes/      <- NoteCard, NoteForm, NoteList
    auth/       <- LoginButton, ProtectedRoute, UserMenu
  pages/
    LoginPage.jsx
    DashboardPage.jsx
    NoteDetailPage.jsx
    NotFoundPage.jsx
  hooks/
    useAuth.js     <- Hook per autenticazione
    useNotes.js    <- Hook per CRUD note
    useApi.js      <- Hook per chiamate API
  context/
    AuthContext.jsx
  services/
    api.js         <- Istanza Axios configurata

## Convenzioni React
- Componenti: PascalCase, un file per componente
- Hook: camelCase con prefisso "use"
- Event handler: prefisso "handle"
- Solo functional component con hooks
- MAI fare fetch/axios nei componenti

## Vincoli
- NON usare localStorage per token JWT
- NON hardcodare l'URL del backend. Usa VITE_API_URL
- Ogni pagina deve avere loading state e error state
- Form con validazione client-side prima del submit
\end{lstlisting}

\begin{notebox}
Il file \texttt{SKILL.md} funziona come il \texttt{\_CONTEXT.md} ma è \textbf{specializzato}. Viene caricato dall'IA solo quando il tipo di lavoro lo richiede (Progressive Disclosure). Questo evita di sovraccaricare il contesto.
\end{notebox}

% ---------------------------------------------------------
\section{Il Contesto del Frontend}

\begin{praticabox}[title={Pratica --- Crea \texttt{\_CONTEXT.md} nel frontend}]
\end{praticabox}

\begin{lstlisting}[language={},title={\texttt{\_CONTEXT.md} per Notes Frontend}]
# Progetto: Notes Frontend

## Scopo
Frontend React per l'applicazione Notes. Si connette al
backend Notes API (http://localhost:3000).

## Tecnologie
React 18 + Vite, React Router 6, Tailwind CSS 4, Axios,
React Context per stato globale (auth).

## Flusso di Autenticazione
1. Utente clicca "Login con Google/GitHub"
2. Redirect al backend /api/auth/google
3. Backend gestisce OAuth, restituisce JWT come cookie
4. Frontend chiama /api/auth/me per il profilo
5. AuthContext salva l'utente in memoria
6. Chiamate API successive includono il cookie

## Pagine
/              LoginPage (se non autenticato)
/dashboard     DashboardPage (lista note + creazione)
/notes/:id     NoteDetailPage (dettaglio + modifica)
*              NotFoundPage

## Vincoli
- NON usare localStorage per token JWT
- Axios con withCredentials: true
- Proxy Vite per evitare problemi CORS
\end{lstlisting}

% ---------------------------------------------------------
\section{Generazione del Frontend}

\begin{praticabox}[title={Pratica --- Generare l'applicazione}]
In Copilot Agent Mode:
\begin{lstlisting}[language={}]
Leggi il _CONTEXT.md e il SKILL.md. Implementa l'intero
frontend React.

Ordine:
 1. Configura Tailwind CSS con Vite
 2. Crea src/services/api.js (Axios configurato)
 3. Crea src/context/AuthContext.jsx
 4. Crea componenti UI base (Button, Input, Card)
 5. Crea layout (Header con UserMenu, PageContainer)
 6. Crea LoginPage con bottoni Google e GitHub
 7. Crea ProtectedRoute
 8. Crea src/hooks/useNotes.js
 9. Crea DashboardPage
10. Crea componenti notes (NoteCard, NoteForm)
11. Crea NoteDetailPage
12. Configura React Router in App.jsx
13. Configura proxy Vite
14. Crea .env con VITE_API_URL
\end{lstlisting}
\end{praticabox}

% ---------------------------------------------------------
\section{Configurazione del Proxy Vite}

Per evitare problemi CORS durante lo sviluppo:

\begin{lstlisting}[language=JavaScript,title={\texttt{vite.config.js}}]
import { defineConfig } from 'vite';
import react from '@vitejs/plugin-react';
import tailwindcss from '@tailwindcss/vite';

export default defineConfig({
  plugins: [react(), tailwindcss()],
  server: {
    port: 5173,
    proxy: {
      '/api': {
        target: 'http://localhost:3000',
        changeOrigin: true,
        secure: false
      }
    }
  }
});
\end{lstlisting}

\begin{tipbox}
Con questo proxy, nel frontend puoi chiamare \texttt{/api/notes} invece di \texttt{http://localhost:3000/api/notes}. Vite inoltra automaticamente le richieste al backend.
\end{tipbox}

% ---------------------------------------------------------
\section{Test dell'Applicazione}

\begin{praticabox}[title={Pratica --- Avvio}]
Apri \textbf{due} terminali:

\textbf{Terminale 1 --- Backend:}
\begin{lstlisting}[language=bash]
cd notes-api && npm run dev
\end{lstlisting}

\textbf{Terminale 2 --- Frontend:}
\begin{lstlisting}[language=bash]
cd notes-frontend && npm run dev
\end{lstlisting}

Apri nel browser: \texttt{http://localhost:5173}
\end{praticabox}

\begin{center}
\begin{tabularx}{\textwidth}{l X}
\toprule
\textbf{Funzionalità} & \textbf{Test} \\
\midrule
Login Page & Vedi i bottoni Google e GitHub \\
OAuth Google & Click $\rightarrow$ redirect $\rightarrow$ login $\rightarrow$ dashboard \\
Dashboard & Lista note (vuota all'inizio) \\
Crea nota & Compila form $\rightarrow$ appare nella lista \\
Dettaglio & Click su nota $\rightarrow$ pagina dettaglio \\
Modifica & Modifica titolo/contenuto $\rightarrow$ salva $\rightarrow$ verifica \\
Elimina & Click elimina $\rightarrow$ conferma $\rightarrow$ scompare \\
Logout & Click logout $\rightarrow$ redirect a login \\
Protezione route & Vai a /dashboard senza login $\rightarrow$ redirect \\
User Menu & Nome e avatar dell'utente nel header \\
\bottomrule
\end{tabularx}
\end{center}

\begin{checkpointbox}
Se riesci a fare login, creare una nota, modificarla ed eliminarla, il frontend è completo.
\end{checkpointbox}

% ---------------------------------------------------------
\section{Iterazioni di Design}

\begin{praticabox}[title={Pratica --- Migliorare l'aspetto}]
\begin{lstlisting}[language={}]
Migliora il design:
1. EmptyState quando non ci sono note
2. Animazioni hover sulle card (scale-[1.02])
3. Tag come badge colorati sulle NoteCard
4. Campo ricerca nella dashboard
5. Data di creazione relativa ("2 ore fa", "ieri")
\end{lstlisting}
\end{praticabox}

\begin{praticabox}[title={Pratica --- Responsive design}]
\begin{lstlisting}[language={}]
Verifica che l'interfaccia sia responsive:
1. Griglia note: 1 col mobile, 2 tablet, 3 desktop
2. Header hamburger su mobile
3. Form full-width su mobile
4. Bottoni con touch target >= 44px
\end{lstlisting}
\end{praticabox}

% ---------------------------------------------------------
\section*{Riepilogo}

\begin{center}
\begin{tabularx}{\textwidth}{l X}
\toprule
\textbf{Aspetto} & \textbf{Dettaglio} \\
\midrule
Framework & React~18 + Vite \\
Styling & Tailwind~CSS~4 \\
Auth & OAuth~2.0 via backend proxy \\
State & React Context (AuthContext) \\
Pagine & Login, Dashboard, NoteDetail, 404 \\
Primo SKILL.md & react-frontend-developer \\
Tempo & $\sim$60--90 minuti \\
\bottomrule
\end{tabularx}
\end{center}
